\section*{Regular spreads as projective quadratic structures}
\addcontentsline{toc}{section}{Regular spreads as projective quadratic structures}

Let $q$ be odd, let $\mu\in\mathbb F_q^\times$ be nonsquare, and write
$K=\mathbb F_{q^2}=\mathbb F_q[t]/(t^2-\mu)$.  Regard $V=K^2$ as a
four-dimensional $\mathbb F_q$-space.  If
\[
\Omega(x,y)=x_1y_2-x_2y_1
\]
and $\beta(x,y)$ is the $t$-coefficient of $\Omega(x,y)$, then $\beta$ is a
nondegenerate alternating $\mathbb F_q$-form.  The one-dimensional
$K$-subspaces of $K^2$ are totally isotropic two-dimensional
$\mathbb F_q$-subspaces and form a regular symplectic spread.

Multiplication by $t$ defines $J$ with
\[
J^2=\mu I,
\qquad
\beta(Jx,Jy)=\mu\beta(x,y).
\]
Thus $\sigma=[J]\in PGSp(4,q)$ is a projective involution.  The spread is
exactly the set of $J$-invariant projective lines.  Since
$K^\times/\mathbb F_q^\times\cong C_{q+1}$ has a unique involution, each
regular field-reduction spread carries a unique projective quadratic structure.
For $q\equiv3\pmod4$ one may take $\mu=-1$, recovering the geometric $i$;
for $q\equiv1\pmod4$ a different nonsquare is required.

Restricting the classical Desarguesian-spread stabilizer to symplectic
similitudes gives
\[
C_{PGSp(4,q)}(\sigma)
\cong C_2\times P\Sigma L_2(q^2),
\qquad
\left|C_{PGSp(4,q)}(\sigma)\right|=2q^2(q^4-1).
\]
Therefore the canonical regular-spread orbit has size
\[
\boxed{
\frac{|PGSp(4,q)|}{|C_{PGSp(4,q)}(\sigma)|}
=\frac{q^2(q^2-1)}2.
}
\]
This proves the orbit-size formula for every odd $q$ and recovers the complete
computational values
\[
36,\quad300,\quad1176
\]
at $q=3,5,7$.  It does not promote the all-$q$ rank-three intersection graph,
and it does not classify non-Desarguesian symplectic spreads.

At $q=3$, the established inner spread stabilizer is $S_6$ of order $720$,
so the full stabilizer is
\[
\boxed{C_{PGSp(4,3)}(\sigma)\cong C_2\times S_6.}
\]
The central $C_2=\langle\sigma\rangle$ is silent on the local spread graph;
the visible quotient $S_6$ is the automorphism group of the Kneser local graph
$K(6,2)$ and its Johnson second subconstituent $J(6,3)$.

\subsection*{The shared phase-inversion controller}
The geometric/representation clock has presentation $C_4:C_2\cong D_4$,
while the earlier $\mu_6$ clock has presentation $C_6:C_2\cong D_{12}$.
If the clocks are independent apart from their common outer inverter, their
combined controller is
\[
\Gamma=(C_4\times C_6):C_2,
\qquad |\Gamma|=48,
\]
where the last factor acts by simultaneous inversion.  Exact enumeration gives
\[
Z(\Gamma)\cong C_2^2,
\qquad
[\Gamma,\Gamma]\cong C_6,
\qquad
\Gamma_{\rm ab}\cong C_2^3.
\]
The subgroups $D_4$ and $D_{12}$ intersect in their common inverter and generate
$\Gamma$.  Although $|\Gamma|=48$, it is not the frame stabilizer
$C_2\times S_4$: their center and derived-subgroup orders are respectively
$(4,6)$ and $(2,12)$.

All statements in this insert are finite-geometric or group-theoretic.  The
order-$48$ synthesis assumes clock independence, is not a physical coupling,
and restores none of the withdrawn charge, colour, generation, flux, or
particle identifications.
